\chapter{Caucher Birkar --- Fields Medal 2018}\label{ch2018caucherbirkar:caucher-birkar-fields-medal-2018-a-comprehensive-report-on-his-work-and-impact}

\begin{center}\rule{0.5\linewidth}{0.5pt}\end{center}

\textbf{Prepared:} July 2026 \textbf{Subject:} Caucher Birkar (b. July
1978, Marivan, Kurdistan Province, Iran) \textbf{Position:} Professor,
Yau Mathematical Sciences Center (YMSC) and BIMSA, Tsinghua University;
Professor, University of Cambridge \textbf{Doctoral Advisor:} Ivan
Fesenko (University of Nottingham) \textbf{Fields Medal Citation:}
\emph{``For the proof of the boundedness of Fano varieties and for
contributions to the minimal model program.''}

\begin{center}\rule{0.5\linewidth}{0.5pt}\end{center}

\section{Table of Contents}\label{ch2018caucherbirkar:table-of-contents}

\begin{enumerate}
\def\labelenumi{\arabic{enumi}.}
\tightlist
\item
 Introduction
\item
 Biographical Background
\item
 Algebraic Geometry and Birational Geometry: An Overview
\item
 The Minimal Model Program
\item
 Existence of Minimal Models for Varieties of Log General Type
\item
 Fano Varieties and the Borisov--Alexeev--Borisov Conjecture
\item
 Singularities of Linear Systems
\item
 Boundedness of Complements
\item
 Anti-Pluricanonical Systems on Fano Varieties
\item
 The Proof of the BAB Conjecture
\item
 Generalized Pairs and Log Calabi--Yau Fibrations
\item
 Positive Characteristic Geometry
\item
 Iitaka's Conjecture and Effectivity Results
\item
 The Stolen Medal Incident
\item
 Impact on Algebraic Geometry
\item
 K-Stability and Moduli of Fano Varieties
\item
 Awards and Recognition
\item
 Selected Major Publications
\item
 Ongoing Work and Future Directions
\item
 Conclusion
\item
 References and Further Reading
\end{enumerate}

\begin{center}\rule{0.5\linewidth}{0.5pt}\end{center}

\section{1. Introduction}\label{ch2018caucherbirkar:introduction}

Caucher Birkar was awarded the Fields Medal in 2018 for his proof of the
boundedness of Fano varieties --- a landmark result in birational
geometry that resolved the Borisov--Alexeev--Borisov (BAB) conjecture, a
problem that had resisted solution for over two decades. His work
substantially advanced the Minimal Model Program (MMP), one of the most
ambitious and important research programs in algebraic geometry.

Birkar's journey to the Fields Medal is as notable as his
mathematics. Born to a Kurdish farming family in rural Iran near the
Iraq border, he taught himself mathematics from borrowed books before
fleeing to the United Kingdom as a refugee. Within eighteen years of
arriving in Britain without knowing English, he had won the highest
honor in mathematics.

His contributions span the core of birational geometry: the existence of
minimal models, the boundedness of Fano varieties, the theory of
singularities, and the development of generalized pairs. His work has
enabled major advances in related fields, including the construction of
moduli spaces of Fano varieties via K-stability, and has deep
connections to differential geometry, arithmetic geometry, and
mathematical physics.

This report provides a detailed account of Birkar's important
contributions, the mathematical context in which they arose, the
techniques he developed, and their lasting impact.

\begin{center}\rule{0.5\linewidth}{0.5pt}\end{center}

\section{2. Biographical Background}\label{ch2018caucherbirkar:biographical-background}

\subsection{2.1 Early Life in Iran}\label{ch2018caucherbirkar:early-life-in-iran}

Caucher Birkar was born Fereydoun Derakhshani in July 1978 in Marivan
County, Kurdistan Province, Iran, near the border with Iraq. He was the
third of six children in a Kurdish farming family that relied on
subsistence agriculture. His parents cultivated rice, wheat, and
vegetables, and the family lived a modest existence in a rural village.

Growing up during the Iran--Iraq War (1980--1988), Birkar's region was
closely affected by the conflict. Education opportunities were limited.
His older brother, an engineering student at the University of Tehran,
played a crucial role by bringing mathematics books when he visited.
Birkar became fascinated with mathematics, working through textbooks on
his own. He later said, ``My parents are farmers; my success in math is
very unlikely.''

\subsection{2.2 University Education and Escape from
Iran}\label{ch2018caucherbirkar:university-education-and-escape-from-iran}

Birkar attended the University of Tehran, where he studied mathematics.
As a Kurd and a convert to Christianity, he faced persecution. The
situation became increasingly dangerous, and in 2000 he fled Iran. He
traveled to Turkey and then to the United Kingdom, where he applied for
political asylum.

Upon arriving in the UK, Birkar could not speak English. He adopted the
name ``Caucher Birkar'' --- which means ``migrant mathematician'' in
Kurdish --- to symbolize his new beginning and as a protective measure.
He enrolled at the University of Nottingham, completing a bachelor's
degree and then a PhD under the supervision of Ivan Fesenko in 2004.

\subsection{2.3 Academic Career}\label{ch2018caucherbirkar:academic-career}

After his PhD, Birkar held a postdoctoral research fellowship at the
University of Warwick from 2004 to 2006. In 2006, he joined the
Department of Pure Mathematics and Mathematical Statistics at the
University of Cambridge, where he became a Professor of Mathematics. He
was appointed a Fellow of Trinity Hall.

In 2021, Birkar joined the Yau Mathematical Sciences Center (YMSC) at
Tsinghua University in Beijing as a tenured professor, while also
maintaining his position at Cambridge. In 2022, he also joined the
Beijing Institute of Mathematical Sciences and Applications (BIMSA). In
2025, he was elected as a foreign academician of the Chinese Academy of
Sciences.

\begin{center}\rule{0.5\linewidth}{0.5pt}\end{center}

\section{3. Algebraic Geometry and Birational Geometry: An
Overview}\label{ch2018caucherbirkar:algebraic-geometry-and-birational-geometry-an-overview}

\subsection{3.1 Algebraic Varieties}\label{ch2018caucherbirkar:algebraic-varieties}

Algebraic geometry studies the solution sets of polynomial equations ---
objects called algebraic varieties. For example, the equation
\(x^2 + y^2 = 1\) defines a circle in the plane, while
\(x^2 + y^2 + z^2 = 1\) defines a sphere in 3D space. When the variables
are allowed to be complex numbers, these objects become complex
manifolds with rich geometric structure.

A central goal of algebraic geometry is to classify algebraic varieties
up to \textbf{birational equivalence}. Two varieties are birationally
equivalent if they are isomorphic on dense open subsets --- essentially,
if they agree after ignoring lower-dimensional subsets. This notion
captures the idea that varieties can be ``almost the same'' even if they
differ in small ways.

\subsection{3.2 The Classification
Problem}\label{ch2018caucherbirkar:the-classification-problem}

The birational classification of algebraic varieties proceeds by
dimension: - \textbf{Dimension 1 (curves):} Classified by Riemann in the
19th century. Every curve is birational to a unique smooth projective
curve, classified by its genus. - \textbf{Dimension 2 (surfaces):}
Classified by the Italian school in the early 20th century (Enriques,
Castelnuovo, etc.). - \textbf{Dimension 3 and higher:} The Minimal Model
Program, initiated by Mori in the 1980s, provides a framework for
classification in all dimensions.

\subsection{3.3 The Kodaira Dimension}\label{ch2018caucherbirkar:the-kodaira-dimension}

A fundamental invariant is the \textbf{Kodaira dimension} \(\kappa(X)\),
which measures the growth rate of the pluricanonical sections. Varieties
are divided into three broad classes: - \textbf{General type}
(\(\kappa = \dim X\)): The ``most complex'' varieties, analogous to
curves of genus \(\ge 2\). - \textbf{Calabi--Yau} (\(\kappa = 0\)):
Varieties with trivial canonical bundle, analogous to elliptic curves. -
\textbf{Fano} (\(\kappa = -\infty\)): Varieties with ample
anti-canonical bundle, analogous to the projective line.

The Minimal Model Program seeks to find a ``minimal'' representative in
each birational equivalence class, which is either a minimal model (with
nef canonical divisor) or a Mori fiber space (a Fano fibration).

\begin{center}\rule{0.5\linewidth}{0.5pt}\end{center}

\section{4. The Minimal Model Program}\label{ch2018caucherbirkar:the-minimal-model-program}

\subsection{4.1 Overview}\label{ch2018caucherbirkar:overview}

The Minimal Model Program (MMP), initiated by Shigefumi Mori in the
1980s, proposes to find within each birational class a variety that is
as simple as possible. The program proceeds through a series of
elementary operations:

\begin{enumerate}
\def\labelenumi{\arabic{enumi}.}
\tightlist
\item
 \textbf{Contractions:} Contract a curve along which the canonical
 divisor \(K_X\) is negative.
\item
 \textbf{Flips:} Replace a ``bad'' contraction with a ``good'' one that
 improves the positivity of \(K_X\).
\item
 \textbf{Repetition:} Repeat until either \(K_X\) becomes nef
 (non-negative on all curves) --- giving a minimal model --- or a Mori
 fiber space is obtained.
\end{enumerate}

The MMP is known to work in dimension 2 (classical surface
classification) and dimension 3 (Mori, 1988; Shokurov, 1990s). For
higher dimensions, the program is not fully established, but major
progress has been made.

\subsection{4.2 The Role of
Singularities}\label{ch2018caucherbirkar:the-role-of-singularities}

A key insight of the MMP is that one must allow mild singularities ---
specifically, \textbf{terminal} and \textbf{canonical} singularities ---
to make progress. Even when starting with a smooth variety, the MMP
operations may produce singular varieties, and the theory must
accommodate them.

The classification of singularities according to their severity
(terminal, canonical, log terminal, log canonical) is central to the
MMP. The most useful class is that of \textbf{klt (Kawamata log
terminal) singularities}, which are mild enough to behave well under the
MMP but general enough to include many cases.

\subsection{4.3 The Finite Generation
Conjecture}\label{ch2018caucherbirkar:the-finite-generation-conjecture}

A fundamental question is whether the \textbf{canonical ring}:

\[
R(X, K_X) = \bigoplus_{m \ge 0} H^0(X, \mathcal{O}_X(mK_X))
\]

is finitely generated. Finite generation is essential for the MMP
because it guarantees the existence of the canonical model and the
termination of certain MMP sequences. The finite generation of the
canonical ring for smooth projective varieties was a long-standing
conjecture, finally proved by Birkar, Cascini, Hacon, and McKernan in
2010.

\begin{center}\rule{0.5\linewidth}{0.5pt}\end{center}

\section{5. Existence of Minimal Models for Varieties of Log General
Type}\label{ch2018caucherbirkar:existence-of-minimal-models-for-varieties-of-log-general-type}

\subsection{5.1 The BCHM Breakthrough}\label{ch2018caucherbirkar:the-bchm-breakthrough}

In 2010, Birkar, together with Paolo Cascini, Christopher Hacon, and
James McKernan, published a landmark paper titled ``Existence of minimal
models for varieties of log general type'' in the \emph{Journal of the
American Mathematical Society}.

The paper proved that the Minimal Model Program works for varieties of
log general type --- that is, varieties whose canonical divisor is big.
This case covers a vast range of important varieties, including all
varieties of general type and many others.

\subsection{5.2 Main Results}\label{ch2018caucherbirkar:main-results}

The BCHM paper established several fundamental results:

\begin{enumerate}
\def\labelenumi{\arabic{enumi}.}
\item
 \textbf{Existence of minimal models:} For any smooth projective
 variety \(X\) of log general type (or more generally, for any klt pair
 \((X, \Delta)\) such that \(K_X + \Delta\) is big), a minimal model
 exists.
\item
 \textbf{Finite generation of canonical rings:} The canonical ring
 \(R(X, K_X)\) is finitely generated for any smooth projective variety
 \(X\). This was one of the important open problems in algebraic
 geometry.
\item
 \textbf{Existence of flips:} Flips exist for klt pairs in all
 dimensions. Prior to BCHM, flips were known to exist only up to
 dimension 3.
\item
 \textbf{Log minimal models:} The results hold for log pairs
 \((X, \Delta)\), which are essential for induction on dimension.
\end{enumerate}

\subsection{5.3 Significance}\label{ch2018caucherbirkar:significance}

The BCHM theorem was a substantial result. It proved the existence of
minimal models in complete generality for the first time (for varieties
of log general type), and it established the finite generation of
canonical rings --- a problem that had been open for over 50 years. The
paper introduced several new techniques, including:

\begin{itemize}
\tightlist
\item
 \textbf{Special termination:} A refined analysis of how MMP sequences
 terminate.
\item
 \textbf{Adjunction for flips:} A method for constructing flips using
 the theory of log pairs.
\item
 \textbf{Inductive arguments:} Using the log version of the MMP to
 prove results in higher dimensions by induction.
\end{itemize}

The finite generation of the canonical ring had been conjectured by
Zariski (1940s) and was known in dimension 2 (Zariski, Mumford) and
dimension 3 (Kawamata, 1980s; Mori, 1988). The BCHM result extended this
to all dimensions.

\begin{center}\rule{0.5\linewidth}{0.5pt}\end{center}

\section{6. Fano Varieties and the Borisov--Alexeev--Borisov
Conjecture}\label{ch2018caucherbirkar:fano-varieties-and-the-borisovalexeevborisov-conjecture}

\subsection{6.1 What is a Fano Variety?}\label{ch2018caucherbirkar:what-is-a-fano-variety}

A \textbf{Fano variety} is a projective algebraic variety whose
anti-canonical divisor \(-K_X\) is ample. Fano varieties are the
building blocks of the MMP: they are the ``atomic'' varieties that
cannot be further simplified. They are also of central importance in
differential geometry (they are exactly the varieties that can
potentially admit Kähler--Einstein metrics with positive Ricci
curvature) and in mathematical physics (they appear as the target spaces
of certain sigma models).

\subsection{6.2 The Conjecture}\label{ch2018caucherbirkar:the-conjecture}

The Borisov--Alexeev--Borisov (BAB) conjecture states:

\begin{quote}
\textbf{For any fixed dimension \(n\) and any fixed set
\(D \subset [0,1]\) satisfying the descending chain condition (DCC), the
set of all \(n\)-dimensional Fano varieties with \(\epsilon\)-log
terminal singularities and coefficients in \(D\) is bounded.}
\end{quote}

Here, ``bounded'' means that there exists a projective morphism
\(Z \to T\) of finite type over \(\mathbb{C}\) such that every such Fano
variety appears as a fiber. In practical terms, boundedness means that
these varieties can be parametrized by a finite number of algebraic
families.

A special case --- the important one --- is the boundedness of
\(n\)-dimensional terminal Fano varieties. This was the case that Birkar
solved.

\subsection{6.3 Why Boundedness Matters}\label{ch2018caucherbirkar:why-boundedness-matters}

Boundedness of Fano varieties is fundamental for several reasons:

\begin{itemize}
\tightlist
\item
 \textbf{The MMP:} Fano varieties are the end products of the MMP.
 Knowing they are bounded means that the outputs of the MMP are
 ``finite'' in a precise sense.
\item
 \textbf{Moduli theory:} Boundedness is the first step towards
 constructing moduli spaces of Fano varieties.
\item
 \textbf{K-stability:} The construction of moduli spaces of K-stable
 Fano varieties depends on boundedness.
\item
 \textbf{Birational automorphism groups:} Boundedness implies the
 Jordan property for Cremona groups.
\end{itemize}

\begin{center}\rule{0.5\linewidth}{0.5pt}\end{center}

\section{7. Singularities of Linear
Systems}\label{ch2018caucherbirkar:singularities-of-linear-systems}

\subsection{7.1 The Problem}\label{ch2018caucherbirkar:the-problem}

Before proving the BAB conjecture, Birkar needed to develop powerful new
tools for understanding the singularities of linear systems. In
algebraic geometry, a \textbf{linear system} is a family of divisors
parametrized by a projective space. The singularities of the members of
a linear system --- how they intersect, where they are most singular ---
are crucial for understanding the global geometry of a variety.

\subsection{7.2 Birkar's Contribution}\label{ch2018caucherbirkar:birkars-contribution}

In his 2016 preprint ``Singularities of linear systems and boundedness
of Fano varieties'' (published in the \emph{Annals of Mathematics} in
2021), Birkar proved a series of deep results about the singularities of
linear systems on Fano varieties. The key theorem states:

\begin{quote}
\textbf{For any \(n\)-dimensional Fano variety \(X\) with terminal
singularities and any integer \(m\), the linear system \(|-mK_X|\) has
at most bounded singularities in a precise sense.}
\end{quote}

More precisely, Birkar introduced the notion of \textbf{``boundedness of
complements''} --- the existence of a divisor \(D \sim -mK_X\) such that
\((X, D)\) has log canonical singularities. This is a powerful technical
tool that limits the possible singularities that can appear.

\subsection{7.3 The Connectedness
Theorem}\label{ch2018caucherbirkar:the-connectedness-theorem}

A key ingredient is Birkar's connectedness theorem for non-klt loci,
which generalizes a classical theorem of Shokurov and Kollár. This
theorem controls how the singular locus of a pair \((X, \Delta)\) is
connected, and it is essential for reducing the boundedness problem to
lower-dimensional cases.

\begin{center}\rule{0.5\linewidth}{0.5pt}\end{center}

\section{8. Boundedness of
Complements}\label{ch2018caucherbirkar:boundedness-of-complements}

\subsection{8.1 Shokurov's Conjecture}\label{ch2018caucherbirkar:shokurovs-conjecture}

A \textbf{complement} of a divisor \(D\) on a variety \(X\) is a divisor
\(D'\) such that \(D' \ge D\) and \((X, D')\) is log canonical. The
boundedness of complements conjecture, due to Shokurov, asks whether for
a given Fano variety \(X\) and a fixed integer \(n\), there exists a
bounded set of possible complements.

Birkar solved this conjecture completely in his 2016 work. The theorem
states:

\begin{quote}
\textbf{For fixed dimension \(n\) and a fixed DCC set
\(D \subset [0,1]\), the set of all possible complements on
\(n\)-dimensional Fano varieties with coefficients in \(D\) is bounded.}
\end{quote}

\subsection{8.2 Why Complements Are
Important}\label{ch2018caucherbirkar:why-complements-are-important}

Complements provide a way to ``approximate'' the singularities of a Fano
variety by log canonical pairs. The boundedness of complements means
that:

\begin{enumerate}
\def\labelenumi{\arabic{enumi}.}
\tightlist
\item
 The singularities of Fano varieties are not arbitrarily bad --- they
 are constrained.
\item
 One can find ``good'' divisors that make the pair log canonical, and
 these divisors come from a bounded family.
\item
 This reduces the BAB conjecture to checking a finite number of cases.
\end{enumerate}

The boundedness of complements is now a fundamental tool in birational
geometry, used extensively by other researchers.

\begin{center}\rule{0.5\linewidth}{0.5pt}\end{center}

\section{9. Anti-Pluricanonical Systems on Fano
Varieties}\label{ch2018caucherbirkar:anti-pluricanonical-systems-on-fano-varieties}

\subsection{9.1 Pluricanonical and Anti-Pluricanonical
Systems}\label{ch2018caucherbirkar:pluricanonical-and-anti-pluricanonical-systems}

For a Fano variety \(X\) with ample \(-K_X\), the linear systems
\(|-mK_X|\) for \(m \ge 1\) are called \textbf{anti-pluricanonical
systems}. These systems give embeddings of \(X\) into projective space
(for large enough \(m\)) and are central to the study of Fano varieties.

A fundamental problem is:

\begin{quote}
\textbf{Given a fixed dimension \(n\), does there exist an integer \(m\)
(depending only on \(n\)) such that for every \(n\)-dimensional Fano
variety \(X\) with terminal singularities, the linear system \(|-mK_X|\)
is non-empty and gives a birational map?}
\end{quote}

\subsection{9.2 Birkar's Theorem}\label{ch2018caucherbirkar:birkars-theorem}

In his 2016 paper ``Anti-pluricanonical systems on Fano varieties''
(published in the \emph{Annals of Mathematics} in 2019), Birkar proved:

\begin{quote}
\textbf{For any fixed dimension \(n\), there exists an integer
\(m = m(n)\) such that for every \(n\)-dimensional Fano variety \(X\)
with terminal singularities, the anti-pluricanonical system \(|-mK_X|\)
is non-empty and defines a birational map onto its image.}
\end{quote}

This was a dramatic improvement over previous results. For surfaces, it
was known that \(m = 3\) works (classical result). For 3-folds,
\(m = 4\) was known. Birkar's result gave a uniform bound in all
dimensions, though the explicit \(m(n)\) is large.

\subsection{9.3 Key Techniques}\label{ch2018caucherbirkar:key-techniques}

The proof introduced several new ideas:

\begin{itemize}
\tightlist
\item
 \textbf{Volume bounds:} Using the inequality
 \(\text{vol}(-K_X) \le C(n)\) for terminal Fano varieties.
\item
 \textbf{Multiplicity bounds:} Controlling the multiplicities of
 divisors in anti-pluricanonical systems.
\item
 \textbf{Inductive arguments:} Reducing the problem to lower dimensions
 using the MMP and the theory of log pairs.
\end{itemize}

This result was a crucial step toward the proof of the BAB conjecture,
as it provided effective bounds on the embeddings of Fano varieties.

\begin{center}\rule{0.5\linewidth}{0.5pt}\end{center}

\section{10. The Proof of the BAB
Conjecture}\label{ch2018caucherbirkar:the-proof-of-the-bab-conjecture}

\subsection{10.1 The Main Theorem}\label{ch2018caucherbirkar:the-main-theorem}

Birkar's proof of the BAB conjecture was published in the \emph{Annals
of Mathematics} in 2021. The main theorem states:

\begin{quote}
\textbf{The set of all \(n\)-dimensional Fano varieties with
\(\epsilon\)-log terminal singularities is bounded. In particular, the
set of all \(n\)-dimensional terminal Fano varieties is bounded.}
\end{quote}

\subsection{10.2 Structure of the Proof}\label{ch2018caucherbirkar:structure-of-the-proof}

The proof is long and technically demanding, running over 100 pages. It
proceeds through several steps:

\begin{enumerate}
\def\labelenumi{\arabic{enumi}.}
\item
 \textbf{Reduction to Q-Fano varieties:} By taking a finite cover, one
 can reduce to the case of Q-Fano varieties (where \(-K_X\) is ample
 Q-Cartier).
\item
 \textbf{Boundedness of volumes:} First, one shows that the volume
 \((-K_X)^n\) is bounded above and below for Fano varieties of fixed
 dimension. The upper bound is a classical result
 (Kollár--Miyaoka--Mori); the lower bound is the hard part.
\item
 \textbf{Boundedness of complements:} Using Shokurov's conjecture
 (proved by Birkar), one shows that for any Fano variety, there exists
 a bounded set of possible complements.
\item
 \textbf{Birational boundedness:} Using the anti-pluricanonical systems
 result, one shows that Fano varieties are birationally bounded ---
 i.e., they are birational to varieties in a bounded family.
\item
 \textbf{Removing indeterminacy:} Using the MMP and the theory of log
 pairs, one shows that the birational maps can be resolved, giving
 actual boundedness.
\item
 \textbf{Generalized pairs:} A key technical innovation is the use of
 \textbf{generalized pairs}, which allow one to work with ``virtual''
 boundary divisors that have better properties than ordinary ones.
\end{enumerate}

\subsection{10.3 Generalized Pairs}\label{ch2018caucherbirkar:generalized-pairs}

One of Birkar's important technical innovations is the theory of
\textbf{generalized pairs}. A generalized pair \((X, B + M)\) consists
of a variety \(X\), an ordinary boundary divisor \(B\), and a
\textbf{nef divisor} \(M\) that comes from a higher-dimensional space
via a morphism.

Generalized pairs have better properties than ordinary pairs:

\begin{itemize}
\tightlist
\item
 They behave better under the MMP (flips, contractions).
\item
 They allow for stronger inductive arguments.
\item
 They unify several previous constructions in birational geometry.
\end{itemize}

The theory of generalized pairs, developed extensively by Birkar, is now
a standard tool in birational geometry.

\begin{center}\rule{0.5\linewidth}{0.5pt}\end{center}

\section{11. Generalized Pairs and Log Calabi--Yau
Fibrations}\label{ch2018caucherbirkar:generalized-pairs-and-log-calabiyau-fibrations}

\subsection{11.1 Log Calabi--Yau
Varieties}\label{ch2018caucherbirkar:log-calabiyau-varieties}

A \textbf{log Calabi--Yau pair} \((X, \Delta)\) is a projective pair
such that \(K_X + \Delta \equiv 0\) (numerically trivial). These appear
naturally in the MMP as the outcome of certain minimal model programs,
and they are intimately connected with Fano varieties via fibrations.

\subsection{11.2 Birkar's Results on Log Calabi--Yau
Fibrations}\label{ch2018caucherbirkar:birkars-results-on-log-calabiyau-fibrations}

In a series of papers beginning with ``Log Calabi--Yau fibrations''
(2018), Birkar studied the geometry of fibrations
\(f: (X, \Delta) \to Z\) where the general fiber is a log Calabi--Yau
pair. He proved:

\begin{itemize}
\tightlist
\item
 \textbf{Boundedness results:} The fibers of such fibrations belong to
 bounded families, given appropriate conditions on the coefficients.
\item
 \textbf{Singularity control:} The singularities of the total space and
 the base are controlled by the singularities of the fibers.
\item
 \textbf{DCC for volumes:} The set of possible volumes of log
 Calabi--Yau pairs satisfies the descending chain condition.
\end{itemize}

These results have important applications to the classification of
algebraic varieties and to the theory of moduli spaces.

\subsection{11.3 Applications to Fano
Fibrations}\label{ch2018caucherbirkar:applications-to-fano-fibrations}

A \textbf{Fano fibration} is a morphism \(X \to Z\) whose fibers are
Fano varieties. These are the other possible output of the MMP (along
with minimal models). Birkar's results on log Calabi--Yau fibrations
imply boundedness results for Fano fibrations: the set of all Fano
fibrations of fixed dimension with certain singularity conditions is
bounded.

\begin{center}\rule{0.5\linewidth}{0.5pt}\end{center}

\section{12. Positive Characteristic
Geometry}\label{ch2018caucherbirkar:positive-characteristic-geometry}

\subsection{12.1 The MMP in Positive
Characteristic}\label{ch2018caucherbirkar:the-mmp-in-positive-characteristic}

The Minimal Model Program was originally developed over the complex
numbers (characteristic 0). In positive characteristic, the program is
much harder because many of the standard tools (resolution of
singularities, vanishing theorems) fail or are not known.

Birkar has made important contributions to the MMP in positive
characteristic, particularly for 3-folds in characteristic \(p > 0\).

\subsection{12.2 Existence of Flips and Minimal
Models}\label{ch2018caucherbirkar:existence-of-flips-and-minimal-models}

In his 2016 paper ``Existence of flips and minimal models for 3-folds in
char p'' (published in \emph{Annales Scientifiques de l'ENS}), Birkar
proved:

\begin{itemize}
\tightlist
\item
 The existence of flips for 3-folds in characteristic \(p \ge 7\).
\item
 The existence of minimal models for 3-folds of log general type in
 characteristic \(p \ge 5\).
\end{itemize}

These results, together with work of Hacon and Xu, nearly complete the
Minimal Model Program for 3-folds in positive characteristic.

\subsection{12.3 Mori Fiber Spaces}\label{ch2018caucherbirkar:mori-fiber-spaces}

In joint work with Joe Waldron (``Existence of Mori fibre spaces for
3-folds in char p'', \emph{Advances in Mathematics}, 2017), Birkar
proved the existence of Mori fiber spaces for 3-folds in positive
characteristic, completing the MMP for 3-folds in large enough
characteristic.

\begin{center}\rule{0.5\linewidth}{0.5pt}\end{center}

\section{13. Iitaka's Conjecture and Effectivity
Results}\label{ch2018caucherbirkar:iitakas-conjecture-and-effectivity-results}

\subsection{13.1 Iitaka's Conjecture}\label{ch2018caucherbirkar:iitakas-conjecture}

Iitaka's \(C_{n,m}\) conjecture is a fundamental problem in birational
geometry. It predicts that for a fibration \(f: X \to Y\) with general
fiber \(F\), the Kodaira dimensions satisfy:

\[
\kappa(X) \ge \kappa(F) + \kappa(Y)
\]

This conjecture is known in many cases but remains open in general.

\subsection{13.2 Birkar and Zhang's
Contribution}\label{ch2018caucherbirkar:birkar-and-zhangs-contribution}

In their 2016 IHES paper ``Effectivity of Iitaka fibrations and
pluricanonical systems of polarized pairs,'' Birkar and De-Qi Zhang
proved:

\begin{itemize}
\tightlist
\item
 The Iitaka fibration is \textbf{effective}: there exists an integer
 \(m\) (depending only on the dimension) such that the pluricanonical
 system \(|mK_X|\) gives the Iitaka fibration.
\item
 This effectively solved half of Iitaka's conjecture, reducing it to
 the special case of Kodaira dimension zero.
\end{itemize}

This result was a major advance. The Iitaka fibration had been known to
exist theoretically (by the work of Hacon--McKernan, Takayama, Tsuji),
but the effectivity --- finding an explicit \(m\) that works uniformly
--- was a new and powerful result.

\subsection{\texorpdfstring{13.3 The \(C_{n,m}\) Conjecture for 3-folds
over Finite
Fields}{13.3 The C\_\{n,m\} Conjecture for 3-folds over Finite Fields}}\label{ch2018caucherbirkar:the-c_nm-conjecture-for-3-folds-over-finite-fields}

With Y. Chen and L. Zhang, Birkar proved Iitaka's \(C_{n,m}\) conjecture
for 3-folds over finite fields, extending the reach of the MMP to
arithmetic settings.

\begin{center}\rule{0.5\linewidth}{0.5pt}\end{center}

\section{14. The Stolen Medal Incident}\label{ch2018caucherbirkar:the-stolen-medal-incident}

\subsection{14.1 The Theft}\label{ch2018caucherbirkar:the-theft}

On August 1, 2018, immediately after the Fields Medal ceremony at the
International Congress of Mathematicians in Rio de Janeiro, Birkar's
medal was stolen from a briefcase in a room at the convention center.
The thief also took his wallet and other personal items.

The theft made international headlines. Birkar later said, ``Now I'm
much more famous than I would have been.''

\subsection{14.2 The Replacement}\label{ch2018caucherbirkar:the-replacement}

The International Mathematical Union (IMU) had a spare display medal in
Rio. In a special ceremony at the ICM, Birkar was presented with this
replacement medal. This led to quips that he was the first person to
receive the Fields Medal twice.

The original medal was never recovered.

\subsection{14.3 Impact}\label{ch2018caucherbirkar:impact}

Despite the unfortunate incident, Birkar handled the situation with
grace and humor. He said that the medal was ``just a symbol'' and that
the real honor was the recognition of his work. The incident
inadvertently increased media coverage of the Fields Medal and of
Birkar's notable life story.

\begin{center}\rule{0.5\linewidth}{0.5pt}\end{center}

\section{15. Impact on Algebraic
Geometry}\label{ch2018caucherbirkar:impact-on-algebraic-geometry}

\subsection{15.1 Impact on the Minimal Model
Program}\label{ch2018caucherbirkar:impact-on-the-minimal-model-program}

Birkar's work, particularly the BCHM theorem, transformed the MMP from a
partially completed program into an effective tool that can be applied in
many contexts. The finite generation of canonical rings and the
existence of minimal models for varieties of log general type are now
taken as foundational results.

\subsection{15.2 Impact on the Study of Fano
Varieties}\label{ch2018caucherbirkar:impact-on-the-study-of-fano-varieties}

The boundedness of Fano varieties was a watershed result. Before Birkar,
it was not even clear that the number of deformation types of terminal
Fano varieties was finite. Now, mathematicians know that:

\begin{itemize}
\tightlist
\item
 The set of all terminal Fano varieties of fixed dimension is bounded.
\item
 There are explicit (though large) bounds on their volumes, degrees,
 and other invariants.
\item
 The singularities of Fano varieties are constrained in precise ways.
\end{itemize}

\subsection{15.3 Impact on K-Stability and Moduli
Theory}\label{ch2018caucherbirkar:impact-on-k-stability-and-moduli-theory}

Birkar's boundedness results are essential for the construction of
moduli spaces of Fano varieties via K-stability. The existence of a
projective moduli space of K-polystable Fano varieties, proved by
Chenyang Xu and others, relies crucially on Birkar's boundedness
theorem.

The moduli space of K-stable Fano varieties is now a major research
area, with connections to differential geometry (Kähler--Einstein
metrics), mirror symmetry, and algebraic geometry.

\subsection{15.4 Impact on the Study of Birational Automorphism
Groups}\label{ch2018caucherbirkar:impact-on-the-study-of-birational-automorphism-groups}

The boundedness of Fano varieties implies the \textbf{Jordan property}
for Cremona groups: for any fixed dimension \(n\), there exists a
constant \(J(n)\) such that every finite subgroup of the Cremona group
\(\text{Bir}(\mathbb{P}^n)\) contains an abelian normal subgroup of
index at most \(J(n)\). This was previously known only for
\(\mathbb{P}^2\) (by work of Serre and others).

\subsection{15.5 Impact on Arithmetic
Geometry}\label{ch2018caucherbirkar:impact-on-arithmetic-geometry}

Birkar's results in positive characteristic and his work on Iitaka's
conjecture have applications to arithmetic geometry, particularly to the
study of algebraic varieties over finite fields and to the geometry of
numbers.

\begin{center}\rule{0.5\linewidth}{0.5pt}\end{center}

\section{16. K-Stability and Moduli of Fano
Varieties}\label{ch2018caucherbirkar:k-stability-and-moduli-of-fano-varieties}

\subsection{16.1 The Connection}\label{ch2018caucherbirkar:the-connection}

K-stability is an algebro-geometric stability condition for Fano
varieties that characterizes the existence of Kähler--Einstein metrics.
The construction of moduli spaces of K-stable Fano varieties depends
critically on boundedness results.

\subsection{16.2 Birkar's Contribution}\label{ch2018caucherbirkar:birkars-contribution-1}

Birkar's proof of the BAB conjecture is the key boundedness input for
the construction of K-moduli spaces. The theorem that K-semistable Fano
varieties of fixed dimension with bounded anti-canonical volume form a
bounded family (proved by Jiang, building on Birkar's work) is the first
step in the construction.

The K-moduli space of Fano varieties is now known to exist as a
projective scheme, and it has applications to:

\begin{itemize}
\tightlist
\item
 \textbf{Kähler geometry:} The moduli space parametrizes Fano manifolds
 admitting Kähler--Einstein metrics.
\item
 \textbf{Mirror symmetry:} The moduli space of Fano varieties is
 expected to be related to the moduli space of Landau--Ginzburg models.
\item
 \textbf{Rational points:} Boundedness results are used in the study of
 rational points on Fano varieties.
\end{itemize}

\subsection{16.3 The Fujita--Li Criterion}\label{ch2018caucherbirkar:the-fujitali-criterion}

The reformulation of K-stability by Fujita and Li, in terms of the
``normalized volume'' of valuations, has made K-stability computable for
many examples. This criterion, combined with Birkar's boundedness, has
led to a detailed understanding of which Fano varieties are K-stable in
dimensions 2 and 3.

\begin{center}\rule{0.5\linewidth}{0.5pt}\end{center}

\section{17. Awards and Recognition}\label{ch2018caucherbirkar:awards-and-recognition}

\textbf{2025:} - Foreign Academician, Chinese Academy of Sciences

\textbf{2019:} - Fellow of the Royal Society - Honorary Doctorate,
Salahaddin University, Erbil

\textbf{2018:} - Fields Medal (International Mathematical Union)

\textbf{2016:} - AMS Moore Prize (for the BCHM paper)

\textbf{2010:} - Philip Leverhulme Prize in Mathematics and Statistics -
Prize of the Fondation Sciences Mathematiques de Paris

\textbf{Other:} - Honorary Professor, University of Nottingham - Cecil
King Travel Scholarship, London Mathematical Society

\begin{center}\rule{0.5\linewidth}{0.5pt}\end{center}

\section{18. Selected Major
Publications}\label{ch2018caucherbirkar:selected-major-publications}

\begin{enumerate}
\def\labelenumi{\arabic{enumi}.}
\item
 \textbf{``Existence of minimal models for varieties of log general
 type''} (with P. Cascini, C. Hacon, and J. McKernan). \emph{Journal of
 the American Mathematical Society} 23(2), 405--468, 2010. ---
 \emph{The BCHM breakthrough proving finite generation of canonical
 rings and existence of minimal models.}
\item
 \textbf{``Anti-pluricanonical systems on Fano varieties.''}
 \emph{Annals of Mathematics} 190(2), 345--463, 2019. --- \emph{Proof
 of uniform boundedness of anti-pluricanonical systems on Fano
 varieties.}
\item
 \textbf{``Singularities of linear systems and boundedness of Fano
 varieties.''} \emph{Annals of Mathematics} 193(2), 347--405, 2021. ---
 \emph{Proof of the BAB conjecture (boundedness of Fano varieties) and
 Shokurov's boundedness of complements conjecture.}
\item
 \textbf{``Existence of flips and minimal models for 3-folds in char
 p.''} \emph{Annales Scientifiques de l'ENS} 49(1), 169--212, 2016. ---
 \emph{Positive characteristic MMP for 3-folds.}
\item
 \textbf{``Existence of Mori fibre spaces for 3-folds in char p''}
 (with J. Waldron). \emph{Advances in Mathematics} 313, 62--101, 2017.
 --- \emph{Completion of the MMP for 3-folds in large characteristic.}
\item
 \textbf{``Effectivity of Iitaka fibrations and pluricanonical systems
 of polarized pairs''} (with D.-Q. Zhang). \emph{Publications
 Mathematiques de l'IHES} 123, 283--331, 2016. --- \emph{Effective
 Iitaka fibrations, reducing Iitaka's conjecture to the Kodaira
 dimension zero case.}
\item
 \textbf{``Log Calabi-Yau fibrations.''} arXiv:1811.10709, 2018. ---
 \emph{Foundational results on the geometry of log Calabi-Yau
 fibrations.}
\item
 \textbf{``Generalised pairs in birational geometry.''}
 arXiv:2008.01008, 2020. --- \emph{Systematic development of the theory
 of generalized pairs.}
\item
 \textbf{``Boundedness and volume of generalised pairs.''}
 arXiv:2103.14935, 2021. --- \emph{Further developments on generalized
 pairs and their volumes.}
\item
 \textbf{``Geometry and moduli of polarised varieties.''}
 arXiv:2006.11238, 2020. --- \emph{Connections between boundedness,
 moduli, and generalized pairs.}
\end{enumerate}

\begin{center}\rule{0.5\linewidth}{0.5pt}\end{center}

\section{19. Ongoing Work and Future
Directions}\label{ch2018caucherbirkar:ongoing-work-and-future-directions}

\subsection{19.1 Generalized Pairs}\label{ch2018caucherbirkar:generalized-pairs-1}

Birkar continues to develop the theory of generalized pairs, which has
become a central tool in birational geometry. Recent work focuses on:

\begin{itemize}
\tightlist
\item
 \textbf{Boundedness results:} Extending the theory to prove
 boundedness for wider classes of varieties.
\item
 \textbf{Volume bounds:} Obtaining explicit bounds on the volumes of
 log Calabi--Yau pairs.
\item
 \textbf{DCC properties:} Proving that certain sets of invariants
 satisfy the descending chain condition.
\end{itemize}

\subsection{19.2 Positive Characteristic}\label{ch2018caucherbirkar:positive-characteristic}

Birkar is actively working on extending the MMP to positive
characteristic. The main challenges are:

\begin{itemize}
\tightlist
\item
 Resolution of singularities in positive characteristic (not known in
 general).
\item
 Vanishing theorems in positive characteristic.
\item
 The minimal model program for 4-folds in positive characteristic.
\end{itemize}

\subsection{19.3 Connections to Physics}\label{ch2018caucherbirkar:connections-to-physics}

Recent work with Seung-Joo Lee (2025) applies birational geometry to
supergravity theories, using boundedness results to derive explicit
bounds on the particle spectrum of 6-dimensional supergravity theories.
This represents a new direction connecting Birkar's pure mathematics to
theoretical physics.

\subsection{19.4 Effective Birational
Geometry}\label{ch2018caucherbirkar:effective-birational-geometry}

In his 2025 survey ``Some open problems in effective birational
geometry'' (EMS Surveys), Birkar outlined several fundamental open
problems:

\begin{itemize}
\tightlist
\item
 Finding explicit bounds for the BAB conjecture.
\item
 Effective versions of the existence of minimal models.
\item
 Computational aspects of birational geometry.
\end{itemize}

\subsection{19.5 Positivity, Singularities, and
Boundedness}\label{ch2018caucherbirkar:positivity-singularities-and-boundedness}

His 2025 paper ``Positivity, singularities, and boundedness'' (submitted
to a volume in honor of James McKernan's 60th birthday) explores recent
connections between these three themes, synthesizing developments from
the past decade.

\begin{center}\rule{0.5\linewidth}{0.5pt}\end{center}

\section{20. Conclusion}\label{ch2018caucherbirkar:conclusion}

Caucher Birkar's journey from a Kurdish farming village in Iran to the
Fields Medal is one of the most notable stories in modern
mathematics. His mathematical achievements are equally notable: he
proved the finite generation of canonical rings, established the
existence of minimal models for varieties of log general type, proved
the boundedness of Fano varieties (the BAB conjecture), and solved
Shokurov's boundedness of complements conjecture.

The Fields Medal citation --- ``For the proof of the boundedness of Fano
varieties and for contributions to the minimal model program'' ---
captures the essence of his contribution, but the impact of his work
extends far beyond these specific results. The BCHM theorem is now a
foundational tool in algebraic geometry. The theory of generalized
pairs, which Birkar developed, is a standard technique used by
researchers worldwide. The boundedness of Fano varieties has enabled the
construction of moduli spaces of K-stable Fano varieties, with
connections to differential geometry and mathematical physics.

Birkar's work is characterized by its depth, technical power, and
geometric insight. He has a notable ability to combine the heavy
machinery of the MMP with delicate combinatorial and analytic estimates,
and to find the right level of generality for a problem. His proofs are
technically demanding but conceptually clear, revealing the underlying
structure of the objects he studies.

As of 2026, Birkar continues to push the boundaries of birational
geometry, working on generalized pairs, positive characteristic
geometry, and connections to theoretical physics. His ongoing work
promises to yield further insights into the structure of algebraic
varieties and their classification.

\begin{center}\rule{0.5\linewidth}{0.5pt}\end{center}

\section{21. References and Further
Reading}\label{ch2018caucherbirkar:references-and-further-reading}

\subsection{Official Sources}\label{ch2018caucherbirkar:official-sources}

\begin{itemize}
\tightlist
\item
 IMU Fields Medal 2018 Announcement:
 https://www.mathunion.org/imu-awards/fields-medal/fields-medals-2018
\item
 Fields Medal Citation for Caucher Birkar:
 https://www.mathunion.org/fileadmin/IMU/Prizes/Fields/2018/Birkar-Citation.pdf
\item
 Jackson, A. ``The Work of Caucher Birkar.'' ICM 2018 Laudatio.
 https://www.mathunion.org/fileadmin/IMU/Prizes/Fields/2018/birkar-final.pdf
\end{itemize}

\subsection{Biographical and Overview
Articles}\label{ch2018caucherbirkar:biographical-and-overview-articles}

\begin{itemize}
\tightlist
\item
 Klarreich, E. ``A Refugee with an Unlikely Path Wins the Fields
 Medal.'' \emph{Quanta Magazine}, August 1, 2018.
\item
 Castelvecchi, D. ``Number-theory prodigy among winners of most coveted
 prize in mathematics.'' \emph{Nature}, August 1, 2018.
\item
 Davis, N. and Zhou, N. ``Former refugee among winners of Fields
 medal.'' \emph{The Guardian}, August 1, 2018.
\item
 Phillips, D. ``World's most prestigious maths medal is stolen
 alongside his wallet minutes after professor wins it.'' \emph{The
 Guardian}, August 1, 2018.
\item
 ICM 2018 News: ``Fields Medal winner Caucher Birkar: `Suddenly I have
 wings.'\,''
\end{itemize}

\subsection{Key Technical References}\label{ch2018caucherbirkar:key-technical-references}

\begin{itemize}
\tightlist
\item
 Birkar, C., Cascini, P., Hacon, C., and McKernan, J. ``Existence of
 minimal models for varieties of log general type.'' \emph{JAMS} 23
 (2010), 405--468.
\item
 Birkar, C. ``Anti-pluricanonical systems on Fano varieties.''
 \emph{Annals of Mathematics} 190 (2019), 345--463.
\item
 Birkar, C. ``Singularities of linear systems and boundedness of Fano
 varieties.'' \emph{Annals of Mathematics} 193 (2021), 347--405.
\item
 Birkar, C. ``Existence of flips and minimal models for 3-folds in char
 p.'' \emph{Annales Scientifiques de l'ENS} 49 (2016), 169--212.
\item
 Birkar, C. ``Log Calabi-Yau fibrations.'' arXiv:1811.10709 (2018).
\item
 Birkar, C. ``Generalised pairs in birational geometry.''
 arXiv:2008.01008 (2020).
\item
 Birkar, C. and Zhang, D.-Q. ``Effectivity of Iitaka fibrations and
 pluricanonical systems of polarized pairs.'' \emph{Pub. Math. IHES}
 123 (2016), 283--331.
\item
 Kollar, J. and Mori, S. \emph{Birational Geometry of Algebraic
 Varieties}. Cambridge Tracts in Mathematics 134, 1998.
\item
 Xu, C. ``K-stability of Fano varieties: an algebro-geometric
 approach.'' \emph{EMS Surveys in Mathematical Sciences} 8 (2021),
 265--354.
\end{itemize}

\begin{center}\rule{0.5\linewidth}{0.5pt}\end{center}

\emph{(End of report)}
